\chapter{Theory}
\label{chap:theory}

This chapter establishes the theoretical foundation on which both the artefact described in this thesis and the research that produced it rest. Four bodies of theory are presented in the order in which they are needed by later chapters. Resource scheduling and constraint programming (\Cref{sec:scheduling}) frame the algorithmic problem class, without which the multi-engine algorithm described in \Cref{chap:findings} cannot be assessed against any standard. Human-in-the-loop automation (\Cref{sec:hitl}) frames the interaction model, without which the suggest-and-override architecture and the discussion of automation limits in \Cref{chap:discussion} would appear as stylistic preferences rather than theoretically required positions. The Transport Management System (TMS) landscape (\Cref{sec:tms}) frames the existing tool category Ressursplanlegger enters and explains why traffic coordinators continue to plan manually despite decades of TMS availability; the fit/gap analysis in \Cref{chap:findings} extends the comparison empirically. Design Science Research (\Cref{sec:dsr}) frames the research stance, legitimising a build-and-validate study before \Cref{chap:method} applies its phases. No theory appears in this chapter that is not invoked downstream, and no later claim about the system or the methodology is left ungrounded here.

\section{Resource Scheduling}
\label{sec:scheduling}

Resource scheduling is the problem of assigning a finite set of resources --- people, machines, or vehicles --- to a set of tasks over time, subject to constraints on resource availability, task requirements, and temporal dependencies \parencite{pinedo2016scheduling}. The discipline encompasses several sub-problems that share this structure: nurse rostering allocates clinical staff to hospital shifts; crew scheduling assigns flight crews to legs of a journey; driver scheduling matches operators to vehicle blocks. In each case, the problem reduces to matching resources to tasks under a layered constraint system, with the objective of either feasibility (a valid plan) or optimality (the best plan within a defined cost or quality function).

Ressursplanlegger fits this canonical structure with one notable extension: every assignment requires both a driver and a vehicle, making the problem a multi-resource assignment rather than a single-resource one. The combinatorial state space grows accordingly --- rather than choosing one resource per task, the algorithm must select a feasible pair, with each element constrained independently and jointly. Assignments are tasks with fixed time windows and resource requirements; drivers and vehicles are resources with competence, availability, and capacity constraints; the objective combines maximal coverage of assignments with balanced satisfaction of preference constraints.

Constraints in this domain divide into hard and soft \parencite{rossi2006constraint}. Hard constraints are inviolable: a driver lacking the required licence cannot be assigned, a vehicle outside active status cannot be used, and no resource can be double-booked. Soft constraints encode preferences whose violation reduces solution quality without rendering the plan infeasible: workload should be balanced across drivers, drivers should ideally receive their designated vehicle, and high-priority assignments should be scheduled before lower-priority ones. Each soft constraint carries a configurable weight, allowing the traffic coordinator to express trade-offs --- a soft constraint with a higher weight is more strongly enforced by the optimisation. The full constraint model implemented in Ressursplanlegger maps directly onto this taxonomy and is detailed in \Cref{chap:findings}.

Resource scheduling at realistic fleet sizes is NP-hard; exact methods rapidly become infeasible as instance size grows \parencite{pinedo2016scheduling}. The practical consequence for system design is that no single algorithm dominates the speed--quality space, which motivates a portfolio of solvers tailored to different operating conditions. A neighbouring problem class deserves explicit delimitation: the Vehicle Routing Problem (VRP) addresses the sequencing of geographically distributed deliveries by a fleet \parencite{braekers2016vrp}. VRP and the present problem share a constraint-driven structure but differ in their decision variable. VRP determines the order in which a vehicle visits a set of locations; the present problem determines which driver and vehicle perform an assignment whose start time and location are fixed by the customer order. Sequencing is therefore not a degree of freedom in Ressursplanlegger's algorithm. VRP literature is referenced once, here, to clarify the boundary; it does not return in subsequent chapters.

The portfolio used in Ressursplanlegger spans three solvers that occupy distinct points on the speed--quality trade-off. A greedy heuristic with $\mathcal{O}(n \times m)$ complexity provides instant feedback at the cost of any optimality guarantee. A complete solver based on constraint programming with Boolean satisfiability (CP-SAT) returns near-optimal solutions within a configurable time limit and is well-suited to standard daily plans \parencite{rossi2006constraint}. A metaheuristic engine drawing on tabu search and simulated annealing, two foundational techniques whose modern formulation traces to \textcite{glover1986future}, scales to large or multi-day instances at higher runtime cost. The implementation details, the benchmarks, and the criteria for selecting between solvers belong in \Cref{chap:findings}; what matters here is that an optimal plan, however well it satisfies the constraint model, is not the same as an \emph{acceptable} plan in this domain. That distinction motivates the human-in-the-loop framing developed in the next section.

\section{Human-in-the-Loop Automation}
\label{sec:hitl}

Human-in-the-loop (HITL) automation describes a class of system in which an automated process generates a recommendation or plan, and a human operator reviews, modifies, and approves the output before it takes effect. \textcite{parasuraman2000automation} formalised this class on a ten-level scale running from fully manual operation to fully autonomous action. Levels one to four encompass varieties of decision support in which the human retains every choice; levels seven to ten describe autonomous execution with diminishing human visibility. Ressursplanlegger operates at levels five to six on this scale: the system generates a complete plan and presents it for approval, but the coordinator can modify any element --- reassigning a driver, swapping a vehicle, removing an assignment --- before the plan is committed.

HITL is not a stylistic preference in the present domain; it is a theoretical requirement. Three properties of transport coordination drive this. First, the work environment is unpredictable: weather changes, customer cancellations, sudden sick leave, and equipment failures regularly invalidate prior plans. Second, the assignment task draws heavily on tacit knowledge that is difficult or impossible to formalise --- which drivers cope with winter conditions on which routes, which clients prefer particular operators, which equipment combinations have a history of working together. Third, coordinators operate under accountability: an algorithm that cannot be inspected, queried, or overruled places authority outside the operator who is held responsible for the day's operations.

The design pattern that follows is suggest-and-override. The algorithm absorbs the combinatorial work --- evaluating thousands of feasible permutations against the constraint model --- and presents a coherent candidate plan; the coordinator applies judgment to the cases the model cannot represent. Ressursplanlegger implements this pattern as its core interaction model: every plan that the algorithm produces is editable on a Gantt-style timeline, with conflict and score information surfaced inline so that overrides are informed rather than blind.

Trust calibration is the empirical foundation of the pattern. \textcite{lee2004trust} argue that operator reliance on an automated system tracks the system's perceived purpose, performance, and process, and that reliance is well-calibrated only when these three dimensions are accurately understood. A system whose decisions appear opaque, whose performance bounds are not communicated, or whose process diverges from the operator's mental model produces either disuse, in which sceptical operators reject the system, or misuse, in which trusting operators rely on it inappropriately. Both modes degrade outcomes. The implication for Ressursplanlegger is that override is not merely a fallback for algorithm failure; it is the mechanism by which trust is built and recalibrated as the operator observes how the system handles representative cases.

HITL also imposes an information requirement: the system must expose its reasoning sufficiently for the coordinator to make corrections that improve, rather than degrade, the plan \parencite{lee2004trust}. In Ressursplanlegger this requirement is met by surfacing the score breakdown for each assignment, an inline conflict view showing which constraints are violated, and a dedicated deviation panel. These elements shape the user interface as much as the algorithm itself --- the architecture, the planning timeline, and the deviation views exist because the interaction model demands them, not for ornamental reasons. A theoretically grounded interaction pattern is necessary, but it is not sufficient on its own; the artefact also enters an existing tool landscape, surveyed in the next section.

\section{Transport Management Systems and the Manual-Planning Persistence}
\label{sec:tms}

\textcite{griffis2007tms} conclude that the value proposition of a TMS sits at the boundary between order capture and freight settlement: a TMS is the system of record for what was promised to the customer and what must be billed for once the work is performed. The category centres on order management, carrier and contract handling, freight billing, and operational reporting; the human assignment of resource to task is, with few exceptions, beyond its scope. Modern TMS implementations have absorbed data-driven extensions for visibility and exception handling \parencite{heinbach2022datadriven}, but the cognitive task of deciding which driver and which vehicle to send on which assignment remains, in practice, the responsibility of a human coordinator working outside the optimisation envelope of the system.

Within that category, the Norwegian transport SME market shows uneven penetration. The seven companies interviewed for this thesis split into three patterns: companies operating on a commercial TMS (the Norlog companies on Timpex; Harlem Solutions on Opptur), companies operating on a bespoke in-house system (Ottem and Nordic Crane), and a company with no system at all (Bergen Bulk Transport, where dispatch is conducted by phone and from memory). Two observations follow. First, the manual-planning problem is not exclusively a TMS-shortfall problem; for a non-trivial share of operators, the relevant absence is of any system whatsoever. Second, even where a commercial TMS is in use, the cognitive work of dispatch happens outside it.

\Cref{tab:tms-vs-needs} compares the capabilities supplied by a commercial TMS in this market against the needs that recur across the interview material. The capabilities column is consistent with the TMS literature \parencite{griffis2007tms,heinbach2022datadriven} and with the systems described in the interviews; the needs column is derived from the interview themes and is presented in detail in \Cref{chap:findings}. A fit/gap analysis --- the technique applied in that chapter --- compares system capabilities against documented user needs in order to expose what an existing tool covers and what it does not. The table here previews the structure that the empirical analysis fills in.

\begin{table}[h]
\centering
\caption{Commercial TMS capabilities against traffic coordinator needs.}
\label{tab:tms-vs-needs}
\begin{tabular}{p{0.45\textwidth}p{0.45\textwidth}}
\toprule
\textbf{What commercial TMS provide} & \textbf{What coordinators need but no TMS supplies} \\
\midrule
Order capture and management & Constraint-aware plan generation \\
Customer and contract handling & Unified capacity overview across drivers and vehicles \\
Freight billing and invoicing & Automated conflict and deviation detection \\
Driver-app dispatch and confirmation & Formalisation of tacit assignment knowledge \\
Operational reporting & Sick-leave and disruption replanning support \\
\bottomrule
\end{tabular}
\end{table}

The substantive question is not whether the gap exists, but why it persists. Three structural reasons emerge from the interview material and the surrounding literature. First, no commercial TMS in this market performs plan generation; the assignment cognition has nowhere to go inside the system, regardless of how modern the surrounding software is. Second, ease-of-use and trust barriers raise the cost of any new tool: legacy systems are described by interviewees as extremely slow under concurrent use, with interfaces that have not aged well, which conditions operators to expect that adopting another system will be painful before it is useful. Third, automation in this domain raises a job-protection concern: coordinators expressed wariness toward systems that might either commit decisions on their behalf or render their role redundant, and sceptical respondents conditioned adoption on retaining manual authority \parencite{lee2004trust}. Taken together, these reasons reframe the manual-planning persistence as a structural rather than incidental fact: traffic coordinators plan manually not from preference, but because the available tool category does not include planning, and because the available systems give them little reason to expect that a new one will.

A short analogy makes the structural absence concrete. Standard nurse-rostering software fits hospital wards with relatively stable shift patterns and a homogeneous skill mix; the same software fits home-care agencies poorly, because the unit of work there is short, geographically distributed, and dependent on client-specific competence that the tool cannot represent. The category exists, but the operating pattern it was built for is not the operating pattern of the home-care agency. Norwegian transport SMEs sit in a comparable position with respect to the global TMS category: the category exists, but the operating pattern it was built for is the high-volume long-haul freight operation, not the mixed-fleet, short-horizon, competence-heavy work of the SME coordinator.

None of the existing systems, and none of the structural reasons that perpetuate manual planning, address the space between an assignment existing in the database and a driver being assigned to it. That space is where Ressursplanlegger fits. Building a tool to occupy this gap, and treating that build as research rather than as engineering alone, requires a methodological frame --- Design Science Research, introduced as theory in the next section and applied in \Cref{chap:method}.

\section{Design Science Research}
\label{sec:dsr}

Design Science Research (DSR) is a research methodology in information systems concerned with the creation and evaluation of artefacts --- systems, models, methods, instantiations --- that address practical problems. \textcite{hevner2004design} formalised the framework around three cycles: a relevance cycle anchored in the application domain, a rigour cycle drawing on the existing knowledge base, and a design cycle in which the artefact is built and assessed iteratively. \textcite{peffers2007dsrm} translated this framework into a six-phase process model running from problem identification, through definition of objectives and design and development, to demonstration, evaluation, and communication. Together, the framework and the process model give DSR its characteristic structure of a build that is also an inquiry.

The methodology fits the present project for two reasons. First, the contribution of the thesis is dual: an artefact (Ressursplanlegger) and the design knowledge generated by building, instrumenting, and assessing it against the needs that motivated it. Neither contribution is reducible to the other --- an artefact without inquiry is engineering, and an inquiry without an artefact loses the situated knowledge that comes from confronting an actual problem with an actual system. Second, DSR provides explicit, named phases that can be mapped onto what was done in this project: problem identification through interviews, definition of objectives through requirements, development through iterative implementation, demonstration through scenario walk-throughs, and evaluation through benchmarking and traceability. This mapping gives the methodology chapter a structure against which the build itself can be measured.

A DSR study can pursue either validation or evaluation, and the difference is consequential. \textcite{wieringa2014dsm} distinguishes validation, which predicts how an artefact would behave in the intended context using methods such as benchmarking, expert review, or traceability, from evaluation, which observes how it actually behaves in deployed use. The present thesis validates: the algorithm is benchmarked against synthetic instances, and the implemented requirements are traced against the needs identified in the interviews. The system has not been deployed in production, and no operational performance data has been collected. This boundary is acknowledged here so that subsequent claims about the artefact remain calibrated to the evidence that supports them.

This section presents DSR as theory; the next chapter applies its phases to the present project. The duplication that would otherwise arise between a theory chapter and a methodology chapter is avoided by keeping the abstract framework here and the concrete application --- which interviews, which decisions, which validation activities --- in \Cref{chap:method}.
